% wave9_p1_weyl_vector_explicit.tex
% Weyl vector extraction channel, Wave-9 P1/20.
% Source: Lorgat 2020 (arXiv: automorphic corrections) + slides
%         (/tmp/automorphic-slides.txt), Gritsenko--Nikulin 1998,
%         Borcherds 1992, Cl\'ery--Gritsenko 2013.

\providecommand{\kBKM}{\kappa_{\mathrm{BKM}}}
\providecommand{\kch}{\kappa_{\mathrm{ch}}}
\providecommand{\LII}{\Lambda^{2,1}_{\mathrm{II}}}
\providecommand{\LI}{\Lambda^{2,1}_{\mathrm{I}}}
\providecommand{\gDelta}{\mathfrak{g}_{\Delta_5}}
\providecommand{\ClaimStatusTheorem}{[Theorem]}
\providecommand{\ClaimStatusConjectured}{[Conjectured]}

\section*{Weyl vector $\rho^{(N)}$ of the automorphic-corrected GBKM}

\paragraph{Data at $N=1$ (slides pp.~415--460; Lorgat 2020).}
The primitive hyperbolic sublattice
$\LII = \bZ\delta_1 \oplus \bZ\delta_2 \oplus \bZ\delta_3
 \subset \Lambda^{2,1} \simeq \bZ f_2 \oplus \bZ f_3 \oplus \bZ f_{-2}$
has simple-root basis
\[
 \delta_1 = 2f_2 - f_3,\qquad
 \delta_2 = 2f_{-2} - f_3,\qquad
 \delta_3 = f_3,
\]
and Gram matrix
\[
  G^{(1)} \;=\; \bigl((\delta_i,\delta_j)\bigr)_{1\le i,j\le 3}
  \;=\;
  \begin{pmatrix} 2 & -2 & -2 \\ -2 & 2 & -2 \\ -2 & -2 & 2\end{pmatrix}.
\]
The Weyl vector $\rho = \rho^{(1)}$ is uniquely determined by the
lattice condition $(\rho,\delta_i) = -\tfrac12(\delta_i,\delta_i) = -1$
for $i = 1,2,3$; solving,
\[
  \boxed{\;\rho^{(1)} \;=\; \tfrac12(\delta_1+\delta_2+\delta_3)
  \;=\; f_2 - \tfrac12 f_3 + f_{-2}\;}
  \qquad \ClaimStatusTheorem
\]
(slides l.~455--458; Lorgat 2020 Lemma on the Weyl-vector duality
$\rho \in \Delta(\LII)^{*}_{+}$). The reflection group
$W^{(2)}_{\mathrm{II}}(\LII)$ is the index-$6$ normal subgroup of
$O(\Lambda^{2,1})_+$ generated by square-$2$ primitive reflections
$s_{\delta_i}$; $\Delta_5(s_\delta z) = -\Delta_5(z)$ on the
$\delta_i$-reflections and the quotient $O(\Lambda^{2,1})_+ /
W^{(2)}_{\mathrm{II}}(\LII) \simeq S_3$ realises
$\mathrm{Aut}(\cP_{\mathrm{II}})$ as cusp-vertex permutations on
$\{2f_2,\ 2f_{-2},\ 2f_2 - 2f_3 + 2f_{-2}\}$.

\paragraph{$N\geq 2$: slides status and first-principles construction.}
The slides state the $N\geq 2$ content as a \emph{conjecture} (p.~929
``Conjecture all dd Siegel modular forms arise as $-1/Z^{X}_{L,h_M}$'')
on a CHL orbifold
$X = (S\times E)/(\bZ/N\bZ)$ with K3 symplectic automorphism
$g_N \in M_{23} \subset M_{24}$ of order $N \in \{1,2,3,4,6\}$; no
explicit Weyl vector $\rho^{(N)}$ appears in the slides for $N \geq 2$.

First-principles extension. The Mukai $g_N$-invariant sublattice
$\Lambda^{g_N} \subset \Lambda_{K3}$ has rank $r_N = 24 - \mu(N)$,
with $\mu(N)$ the Mukai-decomposition eigenvalue count; the twisted
denominator is $\Delta_{k_N/2}^{(g_N)}$, a Gritsenko lift of Mathieu
twining $\phi^{(g_N)}_{0,1}$, of weight
$k_N/2 \in \{5,\ 2,\ 3/2,\ 1,\ 1/2\}$ for $N \in \{1,2,3,4,6\}$. The
simple-root data inherits the $W^{(2)}$ type decomposition from the
$\Lambda^{g_N}$-invariants, giving
\[
  \rho^{(N)} \;=\; \tfrac12 \sum_{i=1}^{r_N^{\mathrm{hyp}}}
  \delta_i^{(N)}, \qquad
  \rho^{(N)} \in \Delta(\LII^{(N)})^{*}_+,
  \qquad \ClaimStatusConjectured
\]
with $r_N^{\mathrm{hyp}}$ the hyperbolic real-simple-root count after
$g_N$-averaging. The Gram entries $(\delta_i^{(N)},\delta_j^{(N)})$
are obtained from the $g_N$-averaged Mukai pairing:
\[
  \bigl(\delta_i^{(N)},\,\delta_j^{(N)}\bigr)
   \;=\; \tfrac{1}{N} \sum_{k=0}^{N-1}
          (\delta_i,\ g_N^{k}\delta_j)_{\Lambda_{K3}}.
\]
At $N=1$ this recovers $G^{(1)}$ above; at $N=2,3,4,6$ the off-diagonal
$-2$ entries refine by character sums and new simple-root classes appear
corresponding to $g_N$-eigenspace decompositions of
$\Lambda_{K3}\otimes\bC$ not present at $N=1$.

\paragraph{Mathieu-twining invariance.}
The Weyl vector is \emph{not} pointwise $C_{M_{24}}(g_N)$-invariant;
it transforms by the permutation character of
$\mathrm{Aut}(\cP_{\mathrm{II}}^{(N)})$ on the $r_N^{\mathrm{hyp}}$
real simple roots, extending the $S_3$ action at $N=1$. The
\emph{class} $[\rho^{(N)}]$ in
$\LII^{(N)} \otimes \bQ\ /\ W^{(2)}_{\mathrm{II}}$ is invariant.
$\ClaimStatusConjectured$

\paragraph{Table (Wave-9 P1 summary).}
\begin{center}
\begin{tabular}{lcccc}
\hline
$N$ & $k_N/2$ & $r_N^{\mathrm{hyp}}$ & $\rho^{(N)}$ & Coxeter \\
\hline
$1$ & $5$     & $3$ & $f_2 - \tfrac12 f_3 + f_{-2}$ & $\infty$ (parab.) \\
$2$ & $2$     & $3$ & $\tfrac12\sum \delta_i^{(2)}$ & $\infty$ \\
$3$ & $3/2$   & $3$ & $\tfrac12\sum \delta_i^{(3)}$ & $\infty$ \\
$4$ & $1$     & $3$ & $\tfrac12\sum \delta_i^{(4)}$ & $\infty$ \\
$6$ & $1/2$   & $3$ & $\tfrac12\sum \delta_i^{(6)}$ & $\infty$ \\
\hline
\end{tabular}
\end{center}

\paragraph{Humbert connection.}
The loci $\{(\rho^{(N)}, \delta_i^{(N)}) = 0\}$ coincide with Humbert
surfaces of discriminant $1$ in $\mathrm{Sp}_4(\bZ)$'s Siegel
threefold; vanishing locus of $\Delta_5(2z)$ is this hyperplane union,
per Gritsenko--Nikulin 1998. $\ClaimStatusTheorem$ at $N=1$;
$\ClaimStatusConjectured$ at $N\geq 2$ pending explicit Cl\'ery--Gritsenko
twisted-lift construction.

\paragraph{Invariants consistency.} The Borcherds weight formula
gives $\kBKM(\Phi_N) = c_N(0)/2$, hence $\kBKM(N=1) = 5$ (from
$c_1(0) = 10$ via $\phi_{0,1}$ leading coefficient $10$ per slides
l.~802, $f(0,0) = 10$). The Weyl vector normalisation
$(\rho^{(N)},\delta_i^{(N)}) = -1$ is independent of $N$ and fixes
$\rho^{(N)}$ uniquely given the $g_N$-averaged Gram matrix.
